% TOP_PROBLEM: {"id": 3128, "title": "Switching reconstruction conjecture", "category_id": 3, "category": {"id": 3, "name": "graph_theory", "display_name": "Graph Theory", "description": "Problems involving graphs, networks, and their properties.", "slug": "graph-theory", "order_index": 3, "created_at": "2026-07-31T15:26:25.671Z"}, "status": "open", "problem_number": "OPG-46951", "set_id": 12}
% FIRST_POSTED: 2026-10-01
% ATTEMPT_STATUS: unresolved
% UnsolvedMath ID 3128; source catalogue code OPG-46951
% !TeX program = lualatex
% GPT-6 Astra Ultra research attempt, 2026-10-01; self-contained partial work.
% Completion estimate: no numerical percentage assigned; full problem unresolved.
\documentclass[11pt,a4paper]{article}
\usepackage[margin=25mm]{geometry}
\usepackage{fontspec}
\setmainfont{lmroman10-regular.otf}[BoldFont=lmroman10-bold.otf,
 ItalicFont=lmroman10-italic.otf,BoldItalicFont=lmroman10-bolditalic.otf]
\usepackage{amsmath,amssymb,mathtools}
\usepackage{unicode-math}
\setmathfont{latinmodern-math.otf}
\AtBeginDocument{\let\setminus\smallsetminus}
\usepackage{xurl,hyperref}
\hypersetup{colorlinks=true,urlcolor=blue}
\setcounter{secnumdepth}{0}
\setlength{\parindent}{0pt}
\setlength{\parskip}{5pt}
\setlength{\emergencystretch}{3em}
\begin{document}
\section*{Switching reconstruction conjecture}\label{3128}
{\small UnsolvedMath ID 3128 · OPG-46951 · Graph Theory · OpenGarden}
\subsection{Definitions and mathematical statement}
Let $G=(V,E)$ be a finite simple undirected graph. For a vertex $v\in V$,
its vertex switching $G^{(v)}$ retains the same vertex set, leaves all
adjacencies within $V\setminus\{v\}$ unchanged, and replaces the neighbour
set of $v$ by its complement within $V\setminus\{v\}$. Thus every edge
incident with $v$ becomes a nonedge and every incident nonedge becomes
an edge; no loop is introduced.

With $[H]$ denoting an unlabelled isomorphism class, define the switching
deck to be the multiset
\[
 \mathcal S(G)=\{\!\{[G^{(v)}]:v\in V\}\!\}.
\]
There is one card for each single-vertex switching, and equal isomorphism
classes are kept with multiplicity. The switched vertex is not marked
in its card, nor is there a common labelling supplied across cards.

The conjecture asserts that for any finite simple graphs $G,H$ on
at least five vertices,
\[
 \mathcal S(G)=\mathcal S(H)\quad\Longrightarrow\quad G\cong H.
\]
A card has the same number of vertices as its original graph; this
operation is different from vertex deletion. Neither connectedness nor
any degree condition is assumed.

The deck contains the results of switching once at each individual
vertex, not all graphs obtainable by arbitrary sequences of switchings.
The target is ordinary graph isomorphism, not merely membership in the
same switching-equivalence class. Consequently both the single-switch
convention and card multiplicities are necessary to fix the exact
reconstruction question.
\subsection{Short English statement}
Is every simple graph on at least five vertices determined by the multiset of graphs obtained by complementing all adjacencies at one vertex?
\subsection{Research attempt}
\paragraph{Scope and process.}
GPT-6 Astra Ultra, 1 October 2026. This attempt keeps the catalogue statement above, checks primary literature, and develops a restricted argument with an adversarial gap audit. The proofs below are the mathematical output of this attempt, not a claim of novelty or external certification.

\paragraph{Literature and attack plan.}
Stanley's primary paper [R1] proves switching reconstruction when the order is not divisible by four. Its remaining divisible-by-four cases must not be mistaken for solved by that theorem. We take a separate elementary route: reconstruct the edge count and degree multiset, then identify the switched vertex on a substantial regular subclass. The computations also diagnose exactly why order four is exceptional for edge counting.

\paragraph{Edge and degree information in the deck.}
Let $n=|V(G)|$, $m=|E(G)|$, and let $m_v$ be the number of edges in the card obtained by switching at $v$. Switching deletes its $d(v)$ incident edges and inserts its $n-1-d(v)$ incident nonedges, giving
\[
 m_v=m+n-1-2d(v). \tag{1}
\]
Sum over all $v$ and use $\sum_vd(v)=2m$:
\[
 \sum_vm_v=(n-4)m+n(n-1). \tag{2}
\]
Since $n\ge5$, the deck therefore determines
\[
 m=\frac{\sum_vm_v-n(n-1)}{n-4},
 \qquad
 \{\!\{d(v)\}\!\}
 =\left\{\!\left\{\frac{m+n-1-m_v}{2}\right\}\!\right\}. \tag{3}
\]
These formulas use only card edge counts, not a correspondence between vertices in different cards. In particular, regularity and its degree are recognizable from an unlabelled switching deck.

\paragraph{An elementary regular reconstruction theorem.}
Suppose $G$ is $d$-regular on $n\ge5$ vertices and
\[
 n\notin\{2d,2d+2\}. \tag{4}
\]
In a card $C=G^{(v)}$, the switched vertex has degree $n-1-d$. An original neighbour of $v$ loses its edge to $v$ and has degree $d-1$; an original nonneighbour gains that edge and has degree $d+1$. Condition (4) says exactly that $n-1-d$ differs from both $d-1$ and $d+1$.

It follows that $v$ is the unique vertex of $C$ whose degree is $n-1-d$. Switch at that identifiable vertex again. Every adjacency incident with it is toggled twice, every other adjacency never changes, and the resulting graph is exactly $G$. This completion is well defined on an unlabelled card because any isomorphism preserves the uniquely specified degree.

If some $H$ shares the deck of $G$, formula (3) shows that $H$ is also $d$-regular. The same identification and inverse switch on a common card therefore recover both $G$ and $H$, proving $G\cong H$. This is a reconstruction proof against arbitrary rivals, not merely against rivals assumed regular beforehand.

\paragraph{Boundary examples and what the condition excludes.}
For the empty graph on $n\ge5$ vertices, each card is a star and its degree-$n-1$ centre is identified uniquely. For the complete graph, the switched vertex is isolated and all other vertices retain degree $n-2$, again giving an immediate inverse switch. For a cycle of order eight, $d=2$ and the switched vertex has degree five while other degrees are one or three, so the theorem works even though the order is divisible by four.

The exceptional equations in (4) are genuine collisions of the diagnostic degree: if $n=2d$, then $n-1-d=d-1$; if $n=2d+2$, it equals $d+1$. Several vertices may then pass the same degree test. These collisions invalidate this particular identification rule, but do not produce a pair of nonisomorphic graphs with the same full deck.

At $n=4$, equation (2) loses its dependence on $m$. The empty graph and the four-cycle indeed have the same switching deck, consisting of four copies of $K_{1,3}$. In the cycle, switching a vertex deletes its two incident cycle edges and adds its former opposite edge; the opposite vertex becomes the star centre. This directly verifies the excluded small-order example.

\paragraph{Verification and first remaining gap.}
The reproducible script \texttt{scripts/check\_attempt\_batch\_graphs\_2026\_10\_01.py} verified (1)--(3) on all 1,024 labelled graphs of order five and recovered all 14 labelled regular examples satisfying (4) from each switched card. This is a finite identity and construction check, not a proof for arbitrary orders.

The multiplicity in (2) is essential: summing only distinct card types would give the wrong coefficient. No step confuses switching with deletion, and every card retains its order. Formula (3) recovers a degree multiset but not an adjacency matrix. Identifying degrees in arbitrary irregular cards is therefore insufficient: one must still decide which vertex was switched, or otherwise prove all possible inverse reconstructions isomorphic.

\paragraph{Final status.}
Edge count and degree multiset are reconstructed for all $n\ge5$, and the regular subclass satisfying (4) is fully handled. The arbitrary divisible-by-four cases beyond the cited theorem, including situations where this diagnostic collides, are not settled here. The full conjecture is \textbf{UNRESOLVED} by this attempt.

\subsection{Sources}
\begin{itemize}
\item[S1] ULAM.AI and the UnsolvedMath Contributors, supplied problems.json, record 3128 (OPG-46951), OpenGarden collection. Catalogue identity and source transcription; CC BY 4.0.
\url{https://huggingface.co/datasets/ulamai/UnsolvedMath}.
\item[S2] Open Problem Garden, Switching reconstruction conjecture. Original problem and discussion; the mathematical formulation below is expanded from the supplied source record.
\url{http://www.openproblemgarden.org/op/switching_reconstruction_conjecture}.
\item[R1] R. P. Stanley, \emph{Reconstruction from Vertex-Switching}, Journal of Combinatorial Theory, Series B 38 (1985), 132--138. Primary author-hosted paper checked for Theorem 3.1 and its scope.
\url{https://math.mit.edu/~rstan/pubs/pubfiles/63.pdf}.
\end{itemize}
\end{document}
